\chapter{Conclusion}

This thesis investigated whether the popularity of Reddit posts can be predicted using semantic text embeddings and structured metadata, and whether such predictions can inform a practical posting pipeline. The results showed that Reddit post reception is predictable to a meaningful extent: the tested models were able to capture non-trivial signal related to upvote ratio, indicating that publicly available post data contains useful information about how communities respond to content.

At the same time, the findings also showed clear limits to this predictability. Post reception is not determined by text alone, but is shaped by a broader social process that includes subreddit-specific norms, moderation practices, timing, visibility, and user behavior. This became particularly evident in the live setting, where offline predictive success did not automatically translate into equally strong real-world performance. The study therefore suggests that Reddit popularity can be modeled only partially and under constrained conditions.

% TODO: call to action, warning of feasibility of such work



maybe add in:


Community norms may be more important than general popularity cues

A useful interpretation of your cross-subreddit results is that success may be highly local. What works in one subreddit may fail in another not because the text is objectively worse, but because “good” posts are defined differently across communities.

That supports a broader conclusion:

popularity prediction on Reddit may be better understood as norm prediction within communities than as universal quality prediction across communities

That is a sharp and defensible discussion claim.
